
\documentclass{article}
\usepackage{colortbl}
\usepackage{makecell}
\usepackage{multirow}
\usepackage{supertabular}

\begin{document}

\newcounter{utterance}

\centering \large Interaction Transcript for game `codenames', experiment `opponent\_difficult', episode 3 with gpt{-}5.4{-}azure.
\vspace{24pt}

{ \footnotesize  \setcounter{utterance}{1}
\setlength{\tabcolsep}{0pt}
\begin{supertabular}{c@{$\;$}|p{.15\linewidth}@{}p{.15\linewidth}p{.15\linewidth}p{.15\linewidth}p{.15\linewidth}p{.15\linewidth}}
    \# & $\;$A & \multicolumn{4}{c}{Game Master} & $\;\:$B\\
    \hline

    \theutterance \stepcounter{utterance}  
    & & & \multicolumn{2}{p{0.3\linewidth}}{
        \cellcolor[rgb]{0.9,0.9,0.9}{
            \makecell[{{p{\linewidth}}}]{
                \texttt{\tiny{[GM$|$GM]}}
                {'hidden': {'team': ['διάσ', 'καγκουρό', 'κλέφτης', 'φεγγάρι', 'διάστημα', 'μαστίγιο', 'βερμούδες', 'κούπα', 'ακίδα'], 'innocent': ['αντίχειρας', 'ελεφαντοστό', 'αλλαγή', 'δαυλός', 'αρκούδα', 'σκύλος', 'Τύπος'], 'opponent': ['ραντεβού', 'σαίξπηρ', 'καρδιά', 'ρώμη', 'μπουλόνι', 'ρεύμα', 'διακομιστής', 'ξενοδοχείο'], 'assassin': ['οικόπεδο']}, 'revealed': {'team': {'team': [], 'innocent': [], 'opponent': [], 'assassin': []}, 'opponent': {'team': [], 'innocent': [], 'opponent': [], 'assassin': []}}}
            }
        }
    }
    & & \\ \\

    \theutterance \stepcounter{utterance}  
    & & \multicolumn{4}{p{0.6\linewidth}}{
        \cellcolor[rgb]{0.9,0.9,0.9}{
            \makecell[{{p{\linewidth}}}]{
                \texttt{\tiny{[P1$\langle$GM]}}
                \texttt{Δώστε μία μονολεκτική ένδειξη που σχετίζεται με κάποιες από τις λέξεις της ομάδας σας από τον παρακάτω κατάλογο. Μπορείτε να επιλέξετε σε πόσες και σε ποιες λέξεις αναφέρεται η ένδειξή σας.} \\
\texttt{Ο στόχος είναι να μαντέψει ο συνεργάτης σας όλες τις λέξεις της ομάδας σας με βάση τις ενδείξεις που του δίνετε, πριν η άλλη ομάδα αποκαλύψει όλες τις δικές της. Αν ο συνεργάτης σας μαντέψει μια λέξη του 'αντιπάλου' ή μια λέξη 'αντιπερισπασμού', η σειρά μαντέματός του τελειώνει και η αντίπαλη ομάδα παίρνει τη σειρά της. Αν ο συνεργάτης σας μαντέψει μια λέξη 'θανάτου', το παιχνίδι τελειώνει και η ομάδα σας χάνει αμέσως.} \\
\texttt{Όσο περισσότερες λέξεις σχετίζονται με την ένδειξή σας, τόσο πιο γρήγορα μπορείτε να κερδίσετε το παιχνίδι, αλλά τόσο πιο δύσκολοι μπορεί να είναι οι συσχετισμοί για να τους μαντέψει ο συνεργάτης σας, οπότε επιλέξτε τη λέξη{-}ένδειξή σας με σύνεση. Η λέξη{-}ένδειξη πρέπει να σχετίζεται σημασιολογικά με τις λέξεις{-}στόχους, δεν μπορεί να είναι μία από τις λέξεις των καταλόγων ή να περιέχει μέρη τους.} \\
\\ 
\texttt{Να δίνετε πάντα τη μονολεκτική σας ένδειξη και τον κατάλογο των σχετικών λέξεων{-}στόχων χωρισμένο με κόμματα με την ακόλουθη μορφή και να κάνετε τις απαντήσεις σας όσο το δυνατόν πιο σύντομες, μην περιλαμβάνετε ποτέ άλλο κείμενο από αυτό που απαιτείται σε αυτή τη μορφή:} \\
\\ 
\texttt{ΥΠΟΔΕΙΞΗ: <ΛΕΞΗ>} \\
\texttt{ΣΤΟΧΟΙ: <ΣΤΟΧΟΙ>} \\
\\ 
\texttt{Οι λέξεις της ομάδας σας είναι: διάσ, καγκουρό, κλέφτης, φεγγάρι, διάστημα, μαστίγιο, βερμούδες, κούπα, ακίδα.} \\
\texttt{Οι λέξεις της ομάδας του αντιπάλου σας είναι: ραντεβού, σαίξπηρ, καρδιά, ρώμη, μπουλόνι, ρεύμα, διακομιστής, ξενοδοχείο.} \\
\texttt{Οι λέξεις αντιπερισπασμού είναι: αντίχειρας, ελεφαντοστό, αλλαγή, δαυλός, αρκούδα, σκύλος, Τύπος.} \\
\texttt{Οι λέξεις θανάτου είναι: οικόπεδο.} \\
            }
        }
    }
    & \\ \\

    \theutterance \stepcounter{utterance}  
    & \multicolumn{4}{p{0.6\linewidth}}{
        \cellcolor[rgb]{0.9,0.9,0.9}{
            \makecell[{{p{\linewidth}}}]{
                \texttt{\tiny{[P1$\rangle$GM]}}
                \texttt{ΥΠΟΔΕΙΞΗ: Αστρονομία} \\
\texttt{ΣΤΟΧΟΙ: διάσ, φεγγάρι, διάστημα} \\
            }
        }
    }
    & & \\ \\

    \theutterance \stepcounter{utterance}  
    & & & \multicolumn{2}{p{0.3\linewidth}}{
        \cellcolor[rgb]{0.9,0.9,0.9}{
            \makecell[{{p{\linewidth}}}]{
                \texttt{\tiny{[GM$|$GM]}}
                \texttt{αστρονομία} \\
            }
        }
    }
    & & \\ \\

    \theutterance \stepcounter{utterance}  
    & & & \multicolumn{2}{p{0.3\linewidth}}{
        \cellcolor[rgb]{0.9,0.9,0.9}{
            \makecell[{{p{\linewidth}}}]{
                \texttt{\tiny{[GM$|$GM]}}
                ['διάσ', 'φεγγάρι', 'διάστημα']
            }
        }
    }
    & & \\ \\

    \theutterance \stepcounter{utterance}  
    & & & \multicolumn{2}{p{0.3\linewidth}}{
        \cellcolor[rgb]{0.9,0.9,0.9}{
            \makecell[{{p{\linewidth}}}]{
                \texttt{\tiny{[GM$|$GM]}}
                {'word': 'διάσ', 'assignment': 'team'}
            }
        }
    }
    & & \\ \\

    \theutterance \stepcounter{utterance}  
    & & & \multicolumn{2}{p{0.3\linewidth}}{
        \cellcolor[rgb]{0.9,0.9,0.9}{
            \makecell[{{p{\linewidth}}}]{
                \texttt{\tiny{[GM$|$GM]}}
                {'word': 'φεγγάρι', 'assignment': 'team'}
            }
        }
    }
    & & \\ \\

    \theutterance \stepcounter{utterance}  
    & & & \multicolumn{2}{p{0.3\linewidth}}{
        \cellcolor[rgb]{0.9,0.9,0.9}{
            \makecell[{{p{\linewidth}}}]{
                \texttt{\tiny{[GM$|$GM]}}
                {'word': 'διάστημα', 'assignment': 'team'}
            }
        }
    }
    & & \\ \\

    \theutterance \stepcounter{utterance}  
    & & \multicolumn{4}{p{0.6\linewidth}}{
        \cellcolor[rgb]{0.9,0.9,0.9}{
            \makecell[{{p{\linewidth}}}]{
                \texttt{\tiny{[GM$\rangle$P2]}}
                \texttt{Δώστε μια λίστα χωρισμένη με κόμματα με έως και 3 λέξεις από την παρακάτω λίστα που σχετίζονται καλύτερα ή συνδέονται στενότερα με τη λέξη 'αστρονομία'. Ξεκινήστε πάντα τη λίστα των εικασιών σας με ΔΟΚΙΜΗ:  και μην συμπεριλάβετε κανένα άλλο κείμενο στην απάντησή σας.} \\
\\ 
\texttt{αλλαγή, διάστημα, βερμούδες, μαστίγιο, αντίχειρας, ραντεβού, καρδιά, ξενοδοχείο, ρώμη, καγκουρό, διακομιστής, φεγγάρι, Τύπος, σαίξπηρ, οικόπεδο, σκύλος, ρεύμα, διάσ, κλέφτης, δαυλός, ελεφαντοστό, αρκούδα, κούπα, ακίδα, μπουλόνι} \\
            }
        }
    }
    & \\ \\

    \theutterance \stepcounter{utterance}  
    & & & \multicolumn{4}{p{0.6\linewidth}}{
        \cellcolor[rgb]{0.9,0.9,0.9}{
            \makecell[{{p{\linewidth}}}]{
                \texttt{\tiny{[GM$\langle$P2]}}
                \texttt{ΔΟΚΙΜΗ: διάστημα, φεγγάρι, διάσ} \\
            }
        }
    }
     \\ \\

    \theutterance \stepcounter{utterance}  
    & & & \multicolumn{2}{p{0.3\linewidth}}{
        \cellcolor[rgb]{0.9,0.9,0.9}{
            \makecell[{{p{\linewidth}}}]{
                \texttt{\tiny{[GM$|$GM]}}
                ['διάστημα', 'φεγγάρι', 'διάσ']
            }
        }
    }
    & & \\ \\

    \theutterance \stepcounter{utterance}  
    & & & \multicolumn{2}{p{0.3\linewidth}}{
        \cellcolor[rgb]{0.9,0.9,0.9}{
            \makecell[{{p{\linewidth}}}]{
                \texttt{\tiny{[GM$|$GM]}}
                {'word': 'διάστημα', 'assignment': 'team'}
            }
        }
    }
    & & \\ \\

    \theutterance \stepcounter{utterance}  
    & & & \multicolumn{2}{p{0.3\linewidth}}{
        \cellcolor[rgb]{0.9,0.9,0.9}{
            \makecell[{{p{\linewidth}}}]{
                \texttt{\tiny{[GM$|$GM]}}
                {'word': 'διάστημα', 'assignment': 'team'}
            }
        }
    }
    & & \\ \\

    \theutterance \stepcounter{utterance}  
    & & & \multicolumn{2}{p{0.3\linewidth}}{
        \cellcolor[rgb]{0.9,0.9,0.9}{
            \makecell[{{p{\linewidth}}}]{
                \texttt{\tiny{[GM$|$GM]}}
                {'word': 'φεγγάρι', 'assignment': 'team'}
            }
        }
    }
    & & \\ \\

    \theutterance \stepcounter{utterance}  
    & & & \multicolumn{2}{p{0.3\linewidth}}{
        \cellcolor[rgb]{0.9,0.9,0.9}{
            \makecell[{{p{\linewidth}}}]{
                \texttt{\tiny{[GM$|$GM]}}
                {'word': 'φεγγάρι', 'assignment': 'team'}
            }
        }
    }
    & & \\ \\

    \theutterance \stepcounter{utterance}  
    & & & \multicolumn{2}{p{0.3\linewidth}}{
        \cellcolor[rgb]{0.9,0.9,0.9}{
            \makecell[{{p{\linewidth}}}]{
                \texttt{\tiny{[GM$|$GM]}}
                {'word': 'διάσ', 'assignment': 'team'}
            }
        }
    }
    & & \\ \\

    \theutterance \stepcounter{utterance}  
    & & & \multicolumn{2}{p{0.3\linewidth}}{
        \cellcolor[rgb]{0.9,0.9,0.9}{
            \makecell[{{p{\linewidth}}}]{
                \texttt{\tiny{[GM$|$GM]}}
                {'word': 'διάσ', 'assignment': 'team'}
            }
        }
    }
    & & \\ \\

    \theutterance \stepcounter{utterance}  
    & & & \multicolumn{2}{p{0.3\linewidth}}{
        \cellcolor[rgb]{0.9,0.9,0.9}{
            \makecell[{{p{\linewidth}}}]{
                \texttt{\tiny{[GM$|$GM]}}
                {'word': 'ραντεβού', 'assignment': 'opponent'}
            }
        }
    }
    & & \\ \\

    \theutterance \stepcounter{utterance}  
    & & & \multicolumn{2}{p{0.3\linewidth}}{
        \cellcolor[rgb]{0.9,0.9,0.9}{
            \makecell[{{p{\linewidth}}}]{
                \texttt{\tiny{[GM$|$GM]}}
                {'word': 'μπουλόνι', 'assignment': 'opponent'}
            }
        }
    }
    & & \\ \\

    \theutterance \stepcounter{utterance}  
    & & & \multicolumn{2}{p{0.3\linewidth}}{
        \cellcolor[rgb]{0.9,0.9,0.9}{
            \makecell[{{p{\linewidth}}}]{
                \texttt{\tiny{[GM$|$GM]}}
                {'hidden': {'team': ['καγκουρό', 'κλέφτης', 'μαστίγιο', 'βερμούδες', 'κούπα', 'ακίδα'], 'innocent': ['αντίχειρας', 'ελεφαντοστό', 'αλλαγή', 'δαυλός', 'αρκούδα', 'σκύλος', 'Τύπος'], 'opponent': ['σαίξπηρ', 'καρδιά', 'ρώμη', 'ρεύμα', 'διακομιστής', 'ξενοδοχείο'], 'assassin': ['οικόπεδο']}, 'revealed': {'team': {'team': ['διάστημα', 'φεγγάρι', 'διάσ'], 'innocent': [], 'opponent': [], 'assassin': []}, 'opponent': {'team': [], 'innocent': [], 'opponent': ['ραντεβού', 'μπουλόνι'], 'assassin': []}}}
            }
        }
    }
    & & \\ \\

    \theutterance \stepcounter{utterance}  
    & & \multicolumn{4}{p{0.6\linewidth}}{
        \cellcolor[rgb]{0.9,0.9,0.9}{
            \makecell[{{p{\linewidth}}}]{
                \texttt{\tiny{[P1$\langle$GM]}}
                \texttt{Οι λέξεις διάστημα, φεγγάρι, διάσ μαντεύτηκαν σωστά. Η σειρά του συμπαίκτη σου έληξε εκεί.} \\
\texttt{Τώρα δώστε ένα άλλο στοιχείο που αφορά κάποιες από τις υπόλοιπες λέξεις της ομάδας σας και έναν κατάλογο των σχετικών λέξεων{-}στόχων. Θυμηθείτε να ξεκινήσετε το στοιχείο σας με 'ΥΠΟΔΕΙΞΗ: ', να βάλετε μια νέα γραμμή και να ξεκινήσετε τον κατάλογο των λέξεων{-}στόχων που χωρίζονται με κόμμα με 'ΣΤΟΧΟΙ: '. Σημείωση: ορισμένες λέξεις έχουν αφαιρεθεί από τις λίστες σε σύγκριση με προηγούμενες αιτήσεις.} \\
\\ 
\texttt{Οι υπόλοιπες λέξεις της ομάδας σας είναι: καγκουρό, κλέφτης, μαστίγιο, βερμούδες, κούπα, ακίδα.} \\
\texttt{Οι υπόλοιπες λέξεις για τον αντίπαλό σας είναι: ραντεβού, σαίξπηρ, καρδιά, ρώμη, μπουλόνι, ρεύμα, διακομιστής, ξενοδοχείο.} \\
\texttt{Οι υπόλοιπες λέξεις που αποσπούν την προσοχή είναι: αντίχειρας, ελεφαντοστό, αλλαγή, δαυλός, αρκούδα, σκύλος, Τύπος.} \\
\texttt{Οι υπόλοιπες λέξεις που σκοτώνουν είναι: οικόπεδο.} \\
            }
        }
    }
    & \\ \\

    \theutterance \stepcounter{utterance}  
    & \multicolumn{4}{p{0.6\linewidth}}{
        \cellcolor[rgb]{0.9,0.9,0.9}{
            \makecell[{{p{\linewidth}}}]{
                \texttt{\tiny{[P1$\rangle$GM]}}
                \texttt{ΥΠΟΔΕΙΞΗ: Τσέπη} \\
\texttt{ΣΤΟΧΟΙ: καγκουρό, κλέφτης, βερμούδες} \\
            }
        }
    }
    & & \\ \\

    \theutterance \stepcounter{utterance}  
    & & & \multicolumn{2}{p{0.3\linewidth}}{
        \cellcolor[rgb]{0.9,0.9,0.9}{
            \makecell[{{p{\linewidth}}}]{
                \texttt{\tiny{[GM$|$GM]}}
                \texttt{τσέπη} \\
            }
        }
    }
    & & \\ \\

    \theutterance \stepcounter{utterance}  
    & & & \multicolumn{2}{p{0.3\linewidth}}{
        \cellcolor[rgb]{0.9,0.9,0.9}{
            \makecell[{{p{\linewidth}}}]{
                \texttt{\tiny{[GM$|$GM]}}
                ['καγκουρό', 'κλέφτης', 'βερμούδες']
            }
        }
    }
    & & \\ \\

    \theutterance \stepcounter{utterance}  
    & & & \multicolumn{2}{p{0.3\linewidth}}{
        \cellcolor[rgb]{0.9,0.9,0.9}{
            \makecell[{{p{\linewidth}}}]{
                \texttt{\tiny{[GM$|$GM]}}
                {'word': 'καγκουρό', 'assignment': 'team'}
            }
        }
    }
    & & \\ \\

    \theutterance \stepcounter{utterance}  
    & & & \multicolumn{2}{p{0.3\linewidth}}{
        \cellcolor[rgb]{0.9,0.9,0.9}{
            \makecell[{{p{\linewidth}}}]{
                \texttt{\tiny{[GM$|$GM]}}
                {'word': 'κλέφτης', 'assignment': 'team'}
            }
        }
    }
    & & \\ \\

    \theutterance \stepcounter{utterance}  
    & & & \multicolumn{2}{p{0.3\linewidth}}{
        \cellcolor[rgb]{0.9,0.9,0.9}{
            \makecell[{{p{\linewidth}}}]{
                \texttt{\tiny{[GM$|$GM]}}
                {'word': 'βερμούδες', 'assignment': 'team'}
            }
        }
    }
    & & \\ \\

    \theutterance \stepcounter{utterance}  
    & & \multicolumn{4}{p{0.6\linewidth}}{
        \cellcolor[rgb]{0.9,0.9,0.9}{
            \makecell[{{p{\linewidth}}}]{
                \texttt{\tiny{[GM$\rangle$P2]}}
                \texttt{Οι λέξεις διάστημα, φεγγάρι, διάσ μαντεύτηκαν σωστά. Η σειρά σου έληξε εκεί.} \\
\texttt{Τώρα δώστε έναν άλλο κατάλογο διαχωρισμένο με κόμματα από τουλάχιστον 1 και έως 3 λέξεις από την ακόλουθη λίστα λέξεων που σχετίζονται καλύτερα ή συνδέονται πιο στενά με τη λέξη 'τσέπη'. Θυμηθείτε να ξεκινήσετε την απάντησή σας με ΔΟΚΙΜΗ: . Σημείωση: ορισμένες λέξεις έχουν αφαιρεθεί από τη λίστα σε σύγκριση με προηγούμενα αιτήματα.} \\
\\ 
\texttt{αλλαγή, βερμούδες, μαστίγιο, αντίχειρας, καρδιά, ξενοδοχείο, ρώμη, καγκουρό, διακομιστής, Τύπος, σαίξπηρ, οικόπεδο, σκύλος, ρεύμα, κλέφτης, δαυλός, ελεφαντοστό, αρκούδα, κούπα, ακίδα} \\
            }
        }
    }
    & \\ \\

    \theutterance \stepcounter{utterance}  
    & & & \multicolumn{4}{p{0.6\linewidth}}{
        \cellcolor[rgb]{0.9,0.9,0.9}{
            \makecell[{{p{\linewidth}}}]{
                \texttt{\tiny{[GM$\langle$P2]}}
                \texttt{ΔΟΚΙΜΗ: αλλαγή, βερμούδες, καγκουρό} \\
            }
        }
    }
     \\ \\

    \theutterance \stepcounter{utterance}  
    & & & \multicolumn{2}{p{0.3\linewidth}}{
        \cellcolor[rgb]{0.9,0.9,0.9}{
            \makecell[{{p{\linewidth}}}]{
                \texttt{\tiny{[GM$|$GM]}}
                ['αλλαγή', 'βερμούδες', 'καγκουρό']
            }
        }
    }
    & & \\ \\

    \theutterance \stepcounter{utterance}  
    & & & \multicolumn{2}{p{0.3\linewidth}}{
        \cellcolor[rgb]{0.9,0.9,0.9}{
            \makecell[{{p{\linewidth}}}]{
                \texttt{\tiny{[GM$|$GM]}}
                {'word': 'αλλαγή', 'assignment': 'innocent'}
            }
        }
    }
    & & \\ \\

    \theutterance \stepcounter{utterance}  
    & & & \multicolumn{2}{p{0.3\linewidth}}{
        \cellcolor[rgb]{0.9,0.9,0.9}{
            \makecell[{{p{\linewidth}}}]{
                \texttt{\tiny{[GM$|$GM]}}
                \texttt{αλλαγή} \\
            }
        }
    }
    & & \\ \\

    \theutterance \stepcounter{utterance}  
    & & & \multicolumn{2}{p{0.3\linewidth}}{
        \cellcolor[rgb]{0.9,0.9,0.9}{
            \makecell[{{p{\linewidth}}}]{
                \texttt{\tiny{[GM$|$GM]}}
                {'word': 'ρώμη', 'assignment': 'opponent'}
            }
        }
    }
    & & \\ \\

    \theutterance \stepcounter{utterance}  
    & & & \multicolumn{2}{p{0.3\linewidth}}{
        \cellcolor[rgb]{0.9,0.9,0.9}{
            \makecell[{{p{\linewidth}}}]{
                \texttt{\tiny{[GM$|$GM]}}
                {'word': 'ρεύμα', 'assignment': 'opponent'}
            }
        }
    }
    & & \\ \\

    \theutterance \stepcounter{utterance}  
    & & & \multicolumn{2}{p{0.3\linewidth}}{
        \cellcolor[rgb]{0.9,0.9,0.9}{
            \makecell[{{p{\linewidth}}}]{
                \texttt{\tiny{[GM$|$GM]}}
                {'hidden': {'team': ['καγκουρό', 'κλέφτης', 'μαστίγιο', 'βερμούδες', 'κούπα', 'ακίδα'], 'innocent': ['αντίχειρας', 'ελεφαντοστό', 'δαυλός', 'αρκούδα', 'σκύλος', 'Τύπος'], 'opponent': ['σαίξπηρ', 'καρδιά', 'διακομιστής', 'ξενοδοχείο'], 'assassin': ['οικόπεδο']}, 'revealed': {'team': {'team': ['διάστημα', 'φεγγάρι', 'διάσ'], 'innocent': ['αλλαγή'], 'opponent': [], 'assassin': []}, 'opponent': {'team': [], 'innocent': [], 'opponent': ['ραντεβού', 'μπουλόνι', 'ρώμη', 'ρεύμα'], 'assassin': []}}}
            }
        }
    }
    & & \\ \\

    \theutterance \stepcounter{utterance}  
    & & \multicolumn{4}{p{0.6\linewidth}}{
        \cellcolor[rgb]{0.9,0.9,0.9}{
            \makecell[{{p{\linewidth}}}]{
                \texttt{\tiny{[P1$\langle$GM]}}
                \texttt{Η λέξη αλλαγή μαντεύτηκε, αλλά είναι μια αθώος λέξη. Η σειρά του συμπαίκτη σου έληξε εκεί.} \\
\texttt{Τώρα δώστε ένα άλλο στοιχείο που αφορά κάποιες από τις υπόλοιπες λέξεις της ομάδας σας και έναν κατάλογο των σχετικών λέξεων{-}στόχων. Θυμηθείτε να ξεκινήσετε το στοιχείο σας με 'ΥΠΟΔΕΙΞΗ: ', να βάλετε μια νέα γραμμή και να ξεκινήσετε τον κατάλογο των λέξεων{-}στόχων που χωρίζονται με κόμμα με 'ΣΤΟΧΟΙ: '. Σημείωση: ορισμένες λέξεις έχουν αφαιρεθεί από τις λίστες σε σύγκριση με προηγούμενες αιτήσεις.} \\
\\ 
\texttt{Οι υπόλοιπες λέξεις της ομάδας σας είναι: καγκουρό, κλέφτης, μαστίγιο, βερμούδες, κούπα, ακίδα.} \\
\texttt{Οι υπόλοιπες λέξεις για τον αντίπαλό σας είναι: σαίξπηρ, καρδιά, ρώμη, ρεύμα, διακομιστής, ξενοδοχείο.} \\
\texttt{Οι υπόλοιπες λέξεις που αποσπούν την προσοχή είναι: αντίχειρας, ελεφαντοστό, δαυλός, αρκούδα, σκύλος, Τύπος.} \\
\texttt{Οι υπόλοιπες λέξεις που σκοτώνουν είναι: οικόπεδο.} \\
            }
        }
    }
    & \\ \\

    \theutterance \stepcounter{utterance}  
    & \multicolumn{4}{p{0.6\linewidth}}{
        \cellcolor[rgb]{0.9,0.9,0.9}{
            \makecell[{{p{\linewidth}}}]{
                \texttt{\tiny{[P1$\rangle$GM]}}
                \texttt{ΥΠΟΔΕΙΞΗ: Μάρσιπος} \\
\texttt{ΣΤΟΧΟΙ: καγκουρό} \\
            }
        }
    }
    & & \\ \\

    \theutterance \stepcounter{utterance}  
    & & & \multicolumn{2}{p{0.3\linewidth}}{
        \cellcolor[rgb]{0.9,0.9,0.9}{
            \makecell[{{p{\linewidth}}}]{
                \texttt{\tiny{[GM$|$GM]}}
                \texttt{μάρσιπος} \\
            }
        }
    }
    & & \\ \\

    \theutterance \stepcounter{utterance}  
    & & & \multicolumn{2}{p{0.3\linewidth}}{
        \cellcolor[rgb]{0.9,0.9,0.9}{
            \makecell[{{p{\linewidth}}}]{
                \texttt{\tiny{[GM$|$GM]}}
                ['καγκουρό']
            }
        }
    }
    & & \\ \\

    \theutterance \stepcounter{utterance}  
    & & & \multicolumn{2}{p{0.3\linewidth}}{
        \cellcolor[rgb]{0.9,0.9,0.9}{
            \makecell[{{p{\linewidth}}}]{
                \texttt{\tiny{[GM$|$GM]}}
                {'word': 'καγκουρό', 'assignment': 'team'}
            }
        }
    }
    & & \\ \\

    \theutterance \stepcounter{utterance}  
    & & \multicolumn{4}{p{0.6\linewidth}}{
        \cellcolor[rgb]{0.9,0.9,0.9}{
            \makecell[{{p{\linewidth}}}]{
                \texttt{\tiny{[GM$\rangle$P2]}}
                \texttt{Η λέξη αλλαγή μαντεύτηκε, αλλά είναι μια αθώος λέξη. Η σειρά σου έληξε εκεί.} \\
\texttt{Τώρα δώστε έναν άλλο κατάλογο διαχωρισμένο με κόμματα από τουλάχιστον 1 και έως 1 λέξεις από την ακόλουθη λίστα λέξεων που σχετίζονται καλύτερα ή συνδέονται πιο στενά με τη λέξη 'μάρσιπος'. Θυμηθείτε να ξεκινήσετε την απάντησή σας με ΔΟΚΙΜΗ: . Σημείωση: ορισμένες λέξεις έχουν αφαιρεθεί από τη λίστα σε σύγκριση με προηγούμενα αιτήματα.} \\
\\ 
\texttt{βερμούδες, μαστίγιο, αντίχειρας, καρδιά, ξενοδοχείο, καγκουρό, διακομιστής, Τύπος, σαίξπηρ, οικόπεδο, σκύλος, κλέφτης, δαυλός, ελεφαντοστό, αρκούδα, κούπα, ακίδα} \\
            }
        }
    }
    & \\ \\

    \theutterance \stepcounter{utterance}  
    & & & \multicolumn{4}{p{0.6\linewidth}}{
        \cellcolor[rgb]{0.9,0.9,0.9}{
            \makecell[{{p{\linewidth}}}]{
                \texttt{\tiny{[GM$\langle$P2]}}
                \texttt{ΔΟΚΙΜΗ: καγκουρό} \\
            }
        }
    }
     \\ \\

    \theutterance \stepcounter{utterance}  
    & & & \multicolumn{2}{p{0.3\linewidth}}{
        \cellcolor[rgb]{0.9,0.9,0.9}{
            \makecell[{{p{\linewidth}}}]{
                \texttt{\tiny{[GM$|$GM]}}
                ['καγκουρό']
            }
        }
    }
    & & \\ \\

    \theutterance \stepcounter{utterance}  
    & & & \multicolumn{2}{p{0.3\linewidth}}{
        \cellcolor[rgb]{0.9,0.9,0.9}{
            \makecell[{{p{\linewidth}}}]{
                \texttt{\tiny{[GM$|$GM]}}
                {'word': 'καγκουρό', 'assignment': 'team'}
            }
        }
    }
    & & \\ \\

    \theutterance \stepcounter{utterance}  
    & & & \multicolumn{2}{p{0.3\linewidth}}{
        \cellcolor[rgb]{0.9,0.9,0.9}{
            \makecell[{{p{\linewidth}}}]{
                \texttt{\tiny{[GM$|$GM]}}
                {'word': 'καγκουρό', 'assignment': 'team'}
            }
        }
    }
    & & \\ \\

    \theutterance \stepcounter{utterance}  
    & & & \multicolumn{2}{p{0.3\linewidth}}{
        \cellcolor[rgb]{0.9,0.9,0.9}{
            \makecell[{{p{\linewidth}}}]{
                \texttt{\tiny{[GM$|$GM]}}
                {'word': 'σαίξπηρ', 'assignment': 'opponent'}
            }
        }
    }
    & & \\ \\

    \theutterance \stepcounter{utterance}  
    & & & \multicolumn{2}{p{0.3\linewidth}}{
        \cellcolor[rgb]{0.9,0.9,0.9}{
            \makecell[{{p{\linewidth}}}]{
                \texttt{\tiny{[GM$|$GM]}}
                {'word': 'ξενοδοχείο', 'assignment': 'opponent'}
            }
        }
    }
    & & \\ \\

    \theutterance \stepcounter{utterance}  
    & & & \multicolumn{2}{p{0.3\linewidth}}{
        \cellcolor[rgb]{0.9,0.9,0.9}{
            \makecell[{{p{\linewidth}}}]{
                \texttt{\tiny{[GM$|$GM]}}
                {'hidden': {'team': ['κλέφτης', 'μαστίγιο', 'βερμούδες', 'κούπα', 'ακίδα'], 'innocent': ['αντίχειρας', 'ελεφαντοστό', 'δαυλός', 'αρκούδα', 'σκύλος', 'Τύπος'], 'opponent': ['καρδιά', 'διακομιστής'], 'assassin': ['οικόπεδο']}, 'revealed': {'team': {'team': ['διάστημα', 'φεγγάρι', 'διάσ', 'καγκουρό'], 'innocent': ['αλλαγή'], 'opponent': [], 'assassin': []}, 'opponent': {'team': [], 'innocent': [], 'opponent': ['ραντεβού', 'μπουλόνι', 'ρώμη', 'ρεύμα', 'σαίξπηρ', 'ξενοδοχείο'], 'assassin': []}}}
            }
        }
    }
    & & \\ \\

    \theutterance \stepcounter{utterance}  
    & & \multicolumn{4}{p{0.6\linewidth}}{
        \cellcolor[rgb]{0.9,0.9,0.9}{
            \makecell[{{p{\linewidth}}}]{
                \texttt{\tiny{[P1$\langle$GM]}}
                \texttt{Η λέξη καγκουρό μαντεύτηκε σωστά. Η σειρά του συμπαίκτη σου έληξε εκεί.} \\
\texttt{Τώρα δώστε ένα άλλο στοιχείο που αφορά κάποιες από τις υπόλοιπες λέξεις της ομάδας σας και έναν κατάλογο των σχετικών λέξεων{-}στόχων. Θυμηθείτε να ξεκινήσετε το στοιχείο σας με 'ΥΠΟΔΕΙΞΗ: ', να βάλετε μια νέα γραμμή και να ξεκινήσετε τον κατάλογο των λέξεων{-}στόχων που χωρίζονται με κόμμα με 'ΣΤΟΧΟΙ: '. Σημείωση: ορισμένες λέξεις έχουν αφαιρεθεί από τις λίστες σε σύγκριση με προηγούμενες αιτήσεις.} \\
\\ 
\texttt{Οι υπόλοιπες λέξεις της ομάδας σας είναι: κλέφτης, μαστίγιο, βερμούδες, κούπα, ακίδα.} \\
\texttt{Οι υπόλοιπες λέξεις για τον αντίπαλό σας είναι: σαίξπηρ, καρδιά, διακομιστής, ξενοδοχείο.} \\
\texttt{Οι υπόλοιπες λέξεις που αποσπούν την προσοχή είναι: αντίχειρας, ελεφαντοστό, δαυλός, αρκούδα, σκύλος, Τύπος.} \\
\texttt{Οι υπόλοιπες λέξεις που σκοτώνουν είναι: οικόπεδο.} \\
            }
        }
    }
    & \\ \\

    \theutterance \stepcounter{utterance}  
    & \multicolumn{4}{p{0.6\linewidth}}{
        \cellcolor[rgb]{0.9,0.9,0.9}{
            \makecell[{{p{\linewidth}}}]{
                \texttt{\tiny{[P1$\rangle$GM]}}
                \texttt{ΥΠΟΔΕΙΞΗ: Σορτς} \\
\texttt{ΣΤΟΧΟΙ: βερμούδες} \\
            }
        }
    }
    & & \\ \\

    \theutterance \stepcounter{utterance}  
    & & & \multicolumn{2}{p{0.3\linewidth}}{
        \cellcolor[rgb]{0.9,0.9,0.9}{
            \makecell[{{p{\linewidth}}}]{
                \texttt{\tiny{[GM$|$GM]}}
                \texttt{σορτς} \\
            }
        }
    }
    & & \\ \\

    \theutterance \stepcounter{utterance}  
    & & & \multicolumn{2}{p{0.3\linewidth}}{
        \cellcolor[rgb]{0.9,0.9,0.9}{
            \makecell[{{p{\linewidth}}}]{
                \texttt{\tiny{[GM$|$GM]}}
                ['βερμούδες']
            }
        }
    }
    & & \\ \\

    \theutterance \stepcounter{utterance}  
    & & & \multicolumn{2}{p{0.3\linewidth}}{
        \cellcolor[rgb]{0.9,0.9,0.9}{
            \makecell[{{p{\linewidth}}}]{
                \texttt{\tiny{[GM$|$GM]}}
                {'word': 'βερμούδες', 'assignment': 'team'}
            }
        }
    }
    & & \\ \\

    \theutterance \stepcounter{utterance}  
    & & \multicolumn{4}{p{0.6\linewidth}}{
        \cellcolor[rgb]{0.9,0.9,0.9}{
            \makecell[{{p{\linewidth}}}]{
                \texttt{\tiny{[GM$\rangle$P2]}}
                \texttt{Η λέξη καγκουρό μαντεύτηκε σωστά. Η σειρά σου έληξε εκεί.} \\
\texttt{Τώρα δώστε έναν άλλο κατάλογο διαχωρισμένο με κόμματα από τουλάχιστον 1 και έως 1 λέξεις από την ακόλουθη λίστα λέξεων που σχετίζονται καλύτερα ή συνδέονται πιο στενά με τη λέξη 'σορτς'. Θυμηθείτε να ξεκινήσετε την απάντησή σας με ΔΟΚΙΜΗ: . Σημείωση: ορισμένες λέξεις έχουν αφαιρεθεί από τη λίστα σε σύγκριση με προηγούμενα αιτήματα.} \\
\\ 
\texttt{βερμούδες, μαστίγιο, αντίχειρας, καρδιά, διακομιστής, Τύπος, οικόπεδο, σκύλος, κλέφτης, δαυλός, ελεφαντοστό, αρκούδα, κούπα, ακίδα} \\
            }
        }
    }
    & \\ \\

    \theutterance \stepcounter{utterance}  
    & & & \multicolumn{4}{p{0.6\linewidth}}{
        \cellcolor[rgb]{0.9,0.9,0.9}{
            \makecell[{{p{\linewidth}}}]{
                \texttt{\tiny{[GM$\langle$P2]}}
                \texttt{ΔΟΚΙΜΗ: βερμούδες} \\
            }
        }
    }
     \\ \\

    \theutterance \stepcounter{utterance}  
    & & & \multicolumn{2}{p{0.3\linewidth}}{
        \cellcolor[rgb]{0.9,0.9,0.9}{
            \makecell[{{p{\linewidth}}}]{
                \texttt{\tiny{[GM$|$GM]}}
                ['βερμούδες']
            }
        }
    }
    & & \\ \\

    \theutterance \stepcounter{utterance}  
    & & & \multicolumn{2}{p{0.3\linewidth}}{
        \cellcolor[rgb]{0.9,0.9,0.9}{
            \makecell[{{p{\linewidth}}}]{
                \texttt{\tiny{[GM$|$GM]}}
                {'word': 'βερμούδες', 'assignment': 'team'}
            }
        }
    }
    & & \\ \\

    \theutterance \stepcounter{utterance}  
    & & & \multicolumn{2}{p{0.3\linewidth}}{
        \cellcolor[rgb]{0.9,0.9,0.9}{
            \makecell[{{p{\linewidth}}}]{
                \texttt{\tiny{[GM$|$GM]}}
                {'word': 'βερμούδες', 'assignment': 'team'}
            }
        }
    }
    & & \\ \\

    \theutterance \stepcounter{utterance}  
    & & & \multicolumn{2}{p{0.3\linewidth}}{
        \cellcolor[rgb]{0.9,0.9,0.9}{
            \makecell[{{p{\linewidth}}}]{
                \texttt{\tiny{[GM$|$GM]}}
                {'word': 'διακομιστής', 'assignment': 'opponent'}
            }
        }
    }
    & & \\ \\

    \theutterance \stepcounter{utterance}  
    & & & \multicolumn{2}{p{0.3\linewidth}}{
        \cellcolor[rgb]{0.9,0.9,0.9}{
            \makecell[{{p{\linewidth}}}]{
                \texttt{\tiny{[GM$|$GM]}}
                {'word': 'καρδιά', 'assignment': 'opponent'}
            }
        }
    }
    & & \\ \\

    \theutterance \stepcounter{utterance}  
    & & & \multicolumn{2}{p{0.3\linewidth}}{
        \cellcolor[rgb]{0.9,0.9,0.9}{
            \makecell[{{p{\linewidth}}}]{
                \texttt{\tiny{[GM$|$GM]}}
                {'hidden': {'team': ['κλέφτης', 'μαστίγιο', 'κούπα', 'ακίδα'], 'innocent': ['αντίχειρας', 'ελεφαντοστό', 'δαυλός', 'αρκούδα', 'σκύλος', 'Τύπος'], 'opponent': [], 'assassin': ['οικόπεδο']}, 'revealed': {'team': {'team': ['διάστημα', 'φεγγάρι', 'διάσ', 'καγκουρό', 'βερμούδες'], 'innocent': ['αλλαγή'], 'opponent': [], 'assassin': []}, 'opponent': {'team': [], 'innocent': [], 'opponent': ['ραντεβού', 'μπουλόνι', 'ρώμη', 'ρεύμα', 'σαίξπηρ', 'ξενοδοχείο', 'διακομιστής', 'καρδιά'], 'assassin': []}}}
            }
        }
    }
    & & \\ \\

    \theutterance \stepcounter{utterance}  
    & & \multicolumn{4}{p{0.6\linewidth}}{
        \cellcolor[rgb]{0.9,0.9,0.9}{
            \makecell[{{p{\linewidth}}}]{
                \texttt{\tiny{[P1$\langle$GM]}}
                \texttt{Η λέξη βερμούδες μαντεύτηκε σωστά. Η σειρά του συμπαίκτη σου έληξε εκεί.} \\
\texttt{Τώρα δώστε ένα άλλο στοιχείο που αφορά κάποιες από τις υπόλοιπες λέξεις της ομάδας σας και έναν κατάλογο των σχετικών λέξεων{-}στόχων. Θυμηθείτε να ξεκινήσετε το στοιχείο σας με 'ΥΠΟΔΕΙΞΗ: ', να βάλετε μια νέα γραμμή και να ξεκινήσετε τον κατάλογο των λέξεων{-}στόχων που χωρίζονται με κόμμα με 'ΣΤΟΧΟΙ: '. Σημείωση: ορισμένες λέξεις έχουν αφαιρεθεί από τις λίστες σε σύγκριση με προηγούμενες αιτήσεις.} \\
\\ 
\texttt{Οι υπόλοιπες λέξεις της ομάδας σας είναι: κλέφτης, μαστίγιο, κούπα, ακίδα.} \\
\texttt{Οι υπόλοιπες λέξεις για τον αντίπαλό σας είναι: καρδιά, διακομιστής.} \\
\texttt{Οι υπόλοιπες λέξεις που αποσπούν την προσοχή είναι: αντίχειρας, ελεφαντοστό, δαυλός, αρκούδα, σκύλος, Τύπος.} \\
\texttt{Οι υπόλοιπες λέξεις που σκοτώνουν είναι: οικόπεδο.} \\
            }
        }
    }
    & \\ \\

    \theutterance \stepcounter{utterance}  
    & \multicolumn{4}{p{0.6\linewidth}}{
        \cellcolor[rgb]{0.9,0.9,0.9}{
            \makecell[{{p{\linewidth}}}]{
                \texttt{\tiny{[P1$\rangle$GM]}}
                \texttt{ΥΠΟΔΕΙΞΗ: Κλοπή} \\
\texttt{ΣΤΟΧΟΙ: κλέφτης} \\
            }
        }
    }
    & & \\ \\

    \theutterance \stepcounter{utterance}  
    & & & \multicolumn{2}{p{0.3\linewidth}}{
        \cellcolor[rgb]{0.9,0.9,0.9}{
            \makecell[{{p{\linewidth}}}]{
                \texttt{\tiny{[GM$|$GM]}}
                \texttt{κλοπή} \\
            }
        }
    }
    & & \\ \\

    \theutterance \stepcounter{utterance}  
    & & & \multicolumn{2}{p{0.3\linewidth}}{
        \cellcolor[rgb]{0.9,0.9,0.9}{
            \makecell[{{p{\linewidth}}}]{
                \texttt{\tiny{[GM$|$GM]}}
                ['κλέφτης']
            }
        }
    }
    & & \\ \\

    \theutterance \stepcounter{utterance}  
    & & & \multicolumn{2}{p{0.3\linewidth}}{
        \cellcolor[rgb]{0.9,0.9,0.9}{
            \makecell[{{p{\linewidth}}}]{
                \texttt{\tiny{[GM$|$GM]}}
                {'word': 'κλέφτης', 'assignment': 'team'}
            }
        }
    }
    & & \\ \\

    \theutterance \stepcounter{utterance}  
    & & & \multicolumn{2}{p{0.3\linewidth}}{
        \cellcolor[rgb]{0.9,0.9,0.9}{
            \makecell[{{p{\linewidth}}}]{
                \texttt{\tiny{[GM$|$GM]}}
                \texttt{opponent has won} \\
            }
        }
    }
    & & \\ \\

\end{supertabular}
}

\end{document}
